Correctness feedback can improve the evaluated small text decision models, but its value depends on the feedback budget, the initial policy and the adaptation method. On DecisionBench at 0.8B, one on-policy training pass gains $\gainSmallPkOneDb$ points, reusing one logged sampled verdict per row for eight optimisation passes gains $\gainSmallLogIpsDb$, and eight passes with fresh verdicts gain $\gainSmallPkDb$. The matched early-stopped labelled adapter gains $\gainSmallSupEsDb$. Repeated feedback often reveals labels, and the expected reward gradient is the supervised gradient weighted by the current probability of the correct answer. The logged comparison is consistent with a benefit from additional verdicts, but changes in policy and objective prevent a purely causal attribution to feedback count.

The controls narrow what the gains mean. Prior correction explains much of the improvement on typed-decisions and the public menu; matched labelled fine-tuning outperforms feedback on Laya where the differences are clearest; and a temperature-fitted released model explains much of the calibration gain. Benchmark artifacts, pre-existing training-pool overlap, limited seeds and historical hyper-parameter selection involving test rows restrict generalisation. This study establishes results for simulated feedback on the evaluated text tasks, not for human feedback, image inputs or unrestricted deployments. The next evidence needed is an independent reproduction of the main comparisons and evaluation with one naturally collected verdict per served answer, followed by modality-specific tests if multimodal claims are intended.
